\documentclass[12pt, a4paper, twocolumn]{article}
% --- PACOTES ---
\usepackage[utf8]{inputenc}
\usepackage[T1]{fontenc}
\usepackage[brazil]{babel}
\usepackage{graphicx}
\usepackage{url}
\usepackage{cite}
\usepackage{geometry}
\geometry{top=3cm, bottom=2cm, left=2cm, right=2cm}
\usepackage{tikz}
\usetikzlibrary{shapes.geometric, arrows.meta, positioning, fit, calc}
\usepackage{hyperref}
\hypersetup{breaklinks=true}
% --- CABEÇALHO ---
\title{\textbf{Sistema de Recomendação Nutricional Baseado em Ontologias para Pacientes com Doença Renal Crônica}}
\author{
    Guilherme França Dib de Oliveira Bessa \\
    \small Orientador: Prof. Dr. Cláudio Gottschalg-Duque \\
    \small \textit{Faculdade de Ciência e Tecnologia em Engenharias (FCTE) -- Universidade de Brasília (UnB)}
}
\date{\today}
\begin{document}
\maketitle
% --- RESUMO ---
\begin{abstract}
A Doença Renal Crônica (DRC) afeta mais de 12 milhões de brasileiros e requer um acompanhamento nutricional individualizado e rigoroso para evitar complicações graves associadas ao consumo inadequado dos principais nutrientes que potencializam problemas renais como potássio, sódio e fósforo. No entanto, um número significativo de pacientes consulta a internet quando se trata de aconselhamento dietético, o que os expõe a informações falsas e potencialmente prejudiciais, quase sempre não personalizadas para cada indivíduo. Nesse sentido, o presente trabalho propõe o desenvolvimento de um Sistema de Recomendação Nutricional baseado em Ontologias e técnicas de Inteligência Artificial para orientação alimentar fornecida a pacientes renais de forma segura e personalizada, em conformidade com as diretrizes clínicas. O artigo também estuda aspectos críticos do gerenciamento nutricional da DRC, analisa como o conhecimento especializado pode ser estruturado com ontologias e discute maneiras de melhorar a precisão das recomendações automatizadas, bem como considerações éticas, clínicas e legais, como a conformidade com a Lei Geral de Proteção de Dados Pessoais (LGPD). A abordagem abrange pesquisa aplicada, exploratória e quantitativa, utilizando revisão bibliográfica sistemática e prototipagem em Python. O sistema desenvolvido visa atuar como ferramenta de apoio nutricional ao paciente, fazendo recomendações confiáveis e promovendo a importância dos nutricionistas como principal tomador de decisões e avaliador dessas recomendações geradas.
\end{abstract}

\vspace{0.3cm}
\noindent \textbf{Palavras-chave:} Sistemas de Recomendação. Ontologias. Nutrição Renal. Inteligência Artificial. Saúde Pública. LGPD.
\vspace{0.3cm}

% --- 1. INTRODUÇÃO ---
\section{Introdução}

A Doença Renal Crônica (DRC) é um dos maiores e mais significativos problemas de saúde pública no Brasil, com uma prevalência de mais de 12 milhões de indivíduos, apresentando um espectro de gravidade que vai desde casos manejáveis tratados com diuréticos até eventos graves que requerem hemodiálise ou transplante, os quais são agravados devido à presença de outros quadros clínicos, como diabetes e hipertensão.

Em decorrência da piora dos quadros devido a outras doenças correlacionadas com a alimentação, a dieta é fundamental para o controle da doença, sendo crucial evitar descompensações graves devido ao consumo excessivo de potássio, sódio, fósforo e açúcares. No entanto, apesar de toda a complexidade de estruturação de planos alimentares, uma grande porcentagem de pacientes recorre à internet e às redes sociais como fonte primária de informação, provocada muitas vezes pela falta de dinheiro e pela presença de insegurança alimentar, com pouca garantia de que as informações disponíveis sejam seguras, validadas por profissionais ou individualizadas\footnote{A individualização do plano alimentar é imprescindível na DRC porque os limites de ingestão de potássio, fósforo, proteínas e sódio variam conforme o estágio da doença (G1--G5), a presença de comorbidades e o tipo de terapia renal substitutiva, tornando inviável um padrão dietético único \cite{vasconcelos2021,kdigo2024}.} para o estágio da doença. Essa situação amplifica riscos clínicos e destaca a necessidade de ferramentas inteligentes que forneçam soluções confiáveis e assertivas.

A insegurança alimentar e nutricional --- entendida como a falta de acesso regular e permanente a alimentos em quantidade e qualidade suficientes para uma vida saudável --- é um problema estrutural no Brasil: dados da Pesquisa Nacional por Amostra de Domicílios (PNAD Contínua) indicam que, em 2022, 33,1 milhões de brasileiros se encontravam em situação de fome, e 27,6\% dos domicílios apresentavam algum grau de insegurança alimentar \cite{ibge2023}. Para pacientes com DRC, essa vulnerabilidade é agravada pelo custo elevado de alimentos compatíveis com as restrições renais e pela dificuldade de acesso a orientação nutricional especializada.

O crescimento dos smartphones, da navegação na internet e principalmente das Inteligências Artificiais está exacerbando os fenômenos informais chamados ``Dr. Google'' e ``Dr. IA'': as pessoas procuram diagnósticos e dietas sem mediação profissional. Conforme documentado por Gottschalg-Duque (2024)\cite{posdoc}, essa dinâmica é especialmente perigosa na nefrologia, uma área em que pequenos erros dietéticos podem levar a consequências graves ou fatais. Nesse contexto, torna-se ainda mais importante construir soluções tecnológicas que organizem, padronizem e validem informações nutricionais, permitindo que o paciente receba orientações compatíveis com seu quadro clínico e suas individualidades nutricionais, culturais e financeiras.

Além disso, conforme demonstrado por Mendes et al. (2021), aplicativos de nutrição amplamente utilizados pelo público geral --- como MyFitnessPal (Under Armour, 2005), FatSecret (Secret Industries, 2007) e Lifesum (Lifesum AB, 2013) --- não incorporam restrições clínicas específicas, não diferenciam estágios da DRC e tampouco representam relações semânticas entre alimentos, nutrientes e riscos clínicos \cite{mendes2021}. Estudos recentes apontam que tais aplicativos falham ao atender populações com condições crônicas, pois não são orientados por conhecimento especializado nem por diretrizes clínicas \cite{mendes2021,chen2019}. Isso reforça a necessidade de ferramentas mais seguras, contextualizadas e alinhadas às exigências da nefrologia.

O presente trabalho tem como objetivo o desenvolvimento de um Sistema de Recomendação Nutricional baseado em ontologias e Inteligência Artificial sob uma perspectiva brasileira, considerando inseguranças alimentares como problemas de saúde pública e investigando mecanismos de personalização acessível. Enquanto aplicativos genéricos ignoram restrições clínicas, um sistema especializado deve lidar com níveis de risco, regras nutricionais complexas e estágios da doença, respeitando perfis individuais de pacientes renais.

As ontologias --- definidas como especificações formais e explícitas de uma conceitualização compartilhada \cite{gruber1993} --- constituem uma estratégia fundamental para organizar conhecimento nutricional e clínico de forma semanticamente estruturada, possibilitando o mapeamento de relações como ``alimento -- nutriente -- risco -- estágio da doença''. Trabalhos recentes mostram que ontologias aplicadas à nutrição permitem integrar diretrizes clínicas, dados do paciente e conhecimento especializado, aumentando a precisão das recomendações \cite{zhang2023,mendes2021,mckensy2022}.

A utilização de técnicas de Inteligência Artificial complementa essa abordagem ao ampliar a capacidade do sistema em gerar recomendações personalizadas e incorporar, de maneira confiável, novas informações clínicas, nutricionais e comportamentais. Estudos internacionais reforçam que abordagens híbridas --- ontologias + IA --- têm apresentado resultados consistentes em doenças crônicas \cite{zhang2023,gogebakan2024}.

No contexto brasileiro, a saúde coletiva e o Sistema Único de Saúde (SUS) defendem o uso de tecnologias digitais que ampliem o acesso da população a orientações seguras e personalizadas, respeitando vulnerabilidades sociais e alimentares \cite{krieger2023}. Assim, a adoção de sistemas de recomendação orientados por conhecimento pode contribuir para democratizar informações alimentares confiáveis.

Além dos requisitos técnicos, sistemas digitais aplicados à saúde devem assegurar governança ética e conformidade legal. A Lei Geral de Proteção de Dados Pessoais (LGPD --- Lei n.º 13.709/2018) é a legislação brasileira que regula o tratamento de dados pessoais por pessoas físicas e jurídicas, com o objetivo de proteger os direitos fundamentais de liberdade, privacidade e livre desenvolvimento da personalidade \cite{brasil2018}. A LGPD classifica dados de saúde como \textit{dados pessoais sensíveis} (art. 5º, inciso II), o que implica que seu tratamento exige consentimento qualificado do titular ou enquadramento em hipóteses legais específicas, como a tutela da saúde ou a realização de estudos em saúde pública (art. 11, inciso II) \cite{brasil2018,botelho2021}. Considerando essa classificação, o sistema proposto deve fundamentar-se em princípios de anonimização, transparência e segurança da informação. Dessa forma, sua função consiste em apoiar o processo de tomada de decisão do usuário, sem substituir o julgamento profissional de nutricionistas ou demais especialistas, mantendo imprescindível a avaliação técnica sobre as recomendações geradas.

% --- 2. TRABALHOS RELACIONADOS ---
\section{Trabalhos Relacionados}

Aplicativos de nutrição amplamente utilizados pelo público geral --- como MyFitnessPal (Under Armour, 2005), Yazio (Yazio GmbH, 2015), Lifesum (Lifesum AB, 2013) e FatSecret (Secret Industries, 2007) --- têm ganhado popularidade por fornecerem contagem de calorias e acompanhamento alimentar. Contudo, tais ferramentas apresentam limitações críticas quando aplicadas a pacientes renais, pois não incorporam diretrizes clínicas, não diferenciam estágios da Doença Renal Crônica (DRC) e não oferecem mecanismos semânticos capazes de representar relações entre alimentos, nutrientes, riscos metabólicos e condições clínicas. Mendes et al. (2021) demonstram que aplicativos de nutrição generalistas não atendem adequadamente populações com condições crônicas, justamente por não integrarem conhecimento nutricional especializado em suas recomendações \cite{mendes2021}. De forma complementar, Chen et al. (2019) identificaram que a maioria dos aplicativos nutricionais disponíveis em lojas digitais não possui validação clínica nem acompanhamento por profissionais de saúde \cite{chen2019}.

No campo científico, diversas pesquisas têm explorado o potencial de ontologias para organizar conhecimento em saúde e melhorar processos de recomendação. Gottschalg-Duque et al. (2023), em manuscrito sobre sistemas inteligentes de saúde, argumentam que ontologias permitem estruturar informações complexas e garantir consistência semântica na recomendação \cite{paper_onto}. De modo semelhante, Zhang et al. (2023) desenvolveram um sistema de recomendação nutricional baseado em ontologias para doenças crônicas, evidenciando melhora em precisão, rastreabilidade e transparência do raciocínio nutricional automatizado \cite{zhang2023}.

Mckensy-Sambola et al. (2022) propuseram um motor de recomendação dietética inteiramente baseado em OWL e na Semantic Web Rule Language (SWRL), capaz de inferir automaticamente o perfil nutricional do usuário e associá-lo a planos alimentares compatíveis, com acurácia global de 0,85 em cinco categorias de sobrepeso e obesidade \cite{mckensy2022}. A arquitetura adotada --- ontologia de dietas (ODR) como base semântica, SWRL para inferência e raciocínio automático pelo motor HermiT --- é diretamente análoga à proposta do presente trabalho, substituindo o domínio da obesidade pelo contexto clínico da DRC.

Göğebakan et al. (2024) desenvolveram a ontologia DIAKID especificamente para pacientes com DRC concomitante a diabetes tipo 2, combinando modelagem OWL com regras SWRL extensivas para calcular o estágio da DRC a partir da Taxa de Filtração Glomerular (TFG), identificar interações medicamentosas e ajustar doses de fármacos com potencial de elevar os níveis séricos de potássio \cite{gogebakan2024}. O trabalho demonstra que a modelagem semântica por estágio da doença é viável e precisa, validada com dados de 249 pacientes, e constitui a referência técnica mais próxima da abordagem proposta neste artigo, diferenciando-se pelo foco em recomendações farmacológicas em vez de alimentares.

Em uma revisão sistemática recente conduzida segundo as diretrizes PRISMA e abrangendo estudos de 2014 a 2024, Vitali et al. (2025) identificaram que ferramentas de interoperabilidade semântica --- ontologias, grafos de conhecimento e motores de raciocínio baseados em OWL e SPARQL --- são componentes essenciais para sistemas de recomendação alimentar personalizados e alinhados a condições clínicas individuais \cite{vitali2025}. A revisão reforça que a combinação entre representação formal do conhecimento e Inteligência Artificial Explicável (XAI) potencializa a transparência e a confiabilidade das recomendações geradas, aspecto central para aplicações clínicas como a aqui proposta.

Estudos mais amplos em sistemas de apoio à decisão nutricional também apontam a importância de integrar ontologias com métodos de Inteligência Artificial para obter recomendações personalizadas. Dessa forma, modelos híbridos combinam inferências baseadas em conhecimento, aprendizado de máquina e dados clínicos, permitindo que sistemas acompanhem alterações no quadro do paciente e atualizações em diretrizes de saúde \cite{mendes2021,zhang2023}.

Outro ponto relevante é o contexto brasileiro. Pesquisas em saúde coletiva, como Krieger (2023), reforçam que tecnologias em saúde devem considerar desigualdades socioeconômicas, insegurança alimentar e a centralidade do SUS como rede de cuidado \cite{krieger2023}. A ausência de aplicativos de nutrição voltados especificamente à população renal brasileira reforça um espaço científico e tecnológico pouco explorado.

Portanto, o presente trabalho diferencia-se por combinar três elementos ainda pouco unificados na literatura nacional: (1) uma perspectiva voltada à saúde renal brasileira, (2) uso de ontologias como base para recomendações nutricionais personalizadas com inferência por SWRL, e (3) conformidade explícita com exigências da LGPD, garantindo segurança e ética no tratamento de dados sensíveis.

% --- 3. REFERENCIAL TEÓRICO ---
\section{Referencial Teórico}

    \subsection{A Nutrição na Doença Renal Crônica}

O manejo nutricional da DRC exige controle rigoroso e individualizado de macronutrientes e micronutrientes, uma vez que os rins comprometidos perdem a capacidade de regular adequadamente o equilíbrio eletrolítico e excretar resíduos metabólicos. As principais restrições dietéticas envolvem a limitação de proteínas, potássio, fósforo e sódio, cujos valores de referência variam conforme o estágio da doença --- classificado de G1 a G5 segundo a Taxa de Filtração Glomerular (TFG) nas diretrizes da Kidney Disease: Improving Global Outcomes (KDIGO) --- e a presença de comorbidades como diabetes e hipertensão \cite{vasconcelos2021,kdigo2024}.

A ingestão proteica exemplifica essa complexidade: nos estágios G3 a G5 sem diálise, a KDIGO recomenda ingestão aproximada de 0,8 g/kg/dia para reduzir a produção de ureia e retardar a progressão da doença; já em pacientes em hemodiálise, a necessidade aumenta, pois aminoácidos são perdidos durante o processo dialítico \cite{vasconcelos2021}. De modo análogo, o controle de potássio é crítico para prevenir hipercalemia --- condição que pode causar arritmias cardíacas potencialmente fatais --- e o de fósforo visa evitar a hiperfosfatemia, associada a calcificações vasculares e ósseas. Essas variações tornam inviável a aplicação de um único padrão dietético a todos os pacientes renais, reforçando a necessidade de sistemas capazes de personalizar recomendações com base em parâmetros clínicos individuais \cite{vasconcelos2021}. Portanto, recomendações nutricionais automatizadas devem sempre ser avaliadas por nutricionistas, garantindo respaldo clínico e segurança ao paciente.

    \subsection{Sistemas de Recomendação}

Sistemas de Recomendação (SR) são técnicas computacionais utilizadas para sugerir itens, informações ou ações relevantes a um usuário com base em dados disponíveis ou conhecimento previamente estruturado. Segundo Aggarwal (2016), esses sistemas podem ser classificados em três principais abordagens: a Filtragem Colaborativa, que identifica padrões de preferência entre grupos de usuários; a Filtragem Baseada em Conteúdo, que utiliza características dos itens para gerar recomendações; e os Modelos Híbridos, que combinam múltiplas técnicas para aumentar a precisão \cite{aggarwal}.

No contexto da saúde --- especialmente em aplicações nutricionais para pacientes com Doença Renal Crônica --- modelos exclusivamente colaborativos ou baseados em popularidade são insuficientes, pois não conseguem avaliar as individualidades de cada paciente, suas realidades socioeconômicas e demais variáveis que devem ser consideradas para tratamentos acessíveis e eficientes.

Kulkarni et al. (2023) argumentam que sistemas baseados em conhecimento são mais adequados para cenários onde a precisão e a segurança são essenciais, permitindo incorporar regras, ontologias e parâmetros clínicos diretamente no processo de recomendação \cite{kulkarni}. Dessa forma, abordagens híbridas e orientadas por conhecimento são as mais indicadas para gerar recomendações personalizadas e clinicamente seguras.

    \subsection{Ontologias na Saúde}

As ontologias são especificações formais e explícitas de uma conceitualização compartilhada \cite{gruber1993}, isto é, estruturas que organizam conceitos e relações dentro de um domínio, facilitando a interpretação das informações por meio de sistemas computacionais de maneira consistente. Na saúde, elas contribuem para padronizar o conhecimento clínico e integrar dados diversos, tornando as recomendações mais seguras e explicáveis. Conforme destacado por Gottschalg-Duque et al. (2025), ontologias permitem representar encadeamentos como alimento $\rightarrow$ nutriente $\rightarrow$ risco $\rightarrow$ estágio da doença, possibilitando a aplicação automática de regras nutricionais \cite{gottschalg2025}.

Na nutrição aplicada à Doença Renal Crônica, esse tipo de modelagem auxilia na definição de alimentos permitidos, restritos ou proibidos, de acordo com parâmetros clínicos. Isso favorece a criação de sistemas de recomendação mais precisos, transparentes e alinhados às diretrizes profissionais, superando limitações de métodos baseados apenas em dados ou popularidade.

    \subsection{Ontologias de Referência: ONS e CKDO}

Duas ontologias públicas disponíveis no repositório BioPortal constituem a base semântica do presente trabalho e merecem descrição aprofundada. Ambas são de \textbf{acesso aberto e gratuito}, o que viabiliza a replicação deste estudo sem custos de licenciamento.

A \textbf{Ontology for Nutritional Studies (ONS)} foi desenvolvida no âmbito do projeto europeu ENPADASI, resultado do consenso acadêmico de 51 centros de pesquisa em nove países, e publicada por Vitali et al. (2018) no periódico \textit{Genes \& Nutrition} \cite{ons}. A ONS organiza os conceitos centrais das ciências nutricionais --- dieta, nutriente, alimento e biomarcador --- por meio de classes e relações formalizadas em OWL, reutilizando termos de ontologias consolidadas como BFO, OBI, FoodOn e ChEBI. Sua arquitetura modular permite a incorporação de domínios adicionais, como nutrigenômica, metabolômica e nutricinética, por meio da adição de módulos independentes. A ONS é membro do Open Biological and Biomedical Ontology Foundry (OBO Foundry) e possibilita a descrição padronizada de estudos de intervenção e observacionais em nutrição, viabilizando a integração semântica de dados provenientes de múltiplas fontes. No contexto deste trabalho, a ONS fornece a terminologia formal para representar alimentos e seus componentes nutricionais na ontologia desenvolvida.

A \textbf{Chronic Kidney Disease Ontology (CKDO)}, igualmente disponível no BioPortal sob licença aberta, organiza o conhecimento clínico sobre a DRC de forma computacionalmente interpretável \cite{ckdo}. A CKDO representa os conceitos fundamentais da doença, incluindo os estágios G1 a G5 conforme a classificação pela TFG, os níveis de albuminúria (A1--A3), complicações associadas como anemia e distúrbios minerais, e os critérios diagnósticos adotados internacionalmente. Essa estrutura permite que sistemas computacionais raciocinem sobre o estágio clínico do paciente e apliquem regras diferenciadas de acordo com o grau de comprometimento renal. No presente sistema, a CKDO serve de referência para a modelagem da classe \texttt{EstagiosDRC} e para a associação entre estágios e restrições alimentares específicas, garantindo aderência às diretrizes clínicas vigentes.

Além da ONS e da CKDO, outras ontologias relevantes no domínio biomédico incluem a \textbf{FoodOn} (ontologia aberta de alimentos, integrada ao OBO Foundry) e a \textbf{SNOMED CT} (terminologia clínica internacional, disponível mediante licenciamento gratuito para países membros da IHTSDO, incluindo o Brasil). O presente trabalho optou pela ONS e pela CKDO por sua disponibilidade irrestrita e por sua cobertura direta dos domínios nutricional e nefrológico, respectivamente.

% --- 4. METODOLOGIA ---
\section{Metodologia}

A presente pesquisa caracteriza-se como aplicada, pois visa o desenvolvimento de um protótipo funcional capaz de auxiliar pacientes com Doença Renal Crônica por meio de recomendações nutricionais estruturadas e seguras. Sua abordagem é predominantemente quantitativa, uma vez que envolve a definição de regras formais, modelagem computacional e validação objetiva das recomendações geradas. Quanto aos objetivos, trata-se de uma pesquisa exploratória, ao investigar o uso de tecnologias disruptivas --- como ontologias e técnicas de Inteligência Artificial --- aplicadas ao contexto da nutrição renal.

Os procedimentos adotados envolveram inicialmente uma revisão bibliográfica sistemática sobre nutrição na DRC, sistemas de recomendação e modelagem semântica em saúde. Essa etapa permitiu mapear diretrizes clínicas, identificar lacunas em soluções existentes e estabelecer fundamentos conceituais necessários para estruturar o conhecimento nutricional. Em seguida, procedeu-se à construção da ontologia OntoDRC em OWL~2, com o auxílio do editor Protégé~5.6, organizando relações entre alimentos, nutrientes, riscos clínicos e estágios da doença segundo as diretrizes KDIGO/KDOQI. A base de alimentos foi populada com 52 itens da dieta brasileira extraídos da Tabela Brasileira de Composição de Alimentos (TACO, 4ª edição, NEPA/Unicamp), com valores de potássio, sódio, fósforo, proteína e calorias por 100g. O motor de inferência foi implementado em Python~3 com a biblioteca \texttt{owlready2}, e a interface web foi desenvolvida com o \textit{framework} Flask, permitindo a entrada de dados do paciente e a visualização das recomendações com justificativas semânticas. O protótipo foi testado com perfis sintéticos cobrindo os cinco estágios da DRC (G1 a G5), validando a progressão das restrições conforme a gravidade da doença.

    \subsection{Aspectos Éticos e Conformidade com a LGPD}

Por lidar com dados relacionados à saúde --- classificados como sensíveis pela LGPD (art. 5º, II) ---, o sistema deve seguir rigorosamente a Lei n.º 13.709/2018 \cite{brasil2018}. Conforme Botelho e Camargo (2021), o tratamento de dados de saúde em ambientes digitais exige consentimento qualificado, anonimização e mecanismos de segurança que impeçam a reidentificação do titular \cite{botelho2021}. Todos os dados utilizados neste protótipo são tratados de forma anonimizada, garantindo privacidade e impossibilitando a identificação direta do usuário. Além disso, o sistema atua apenas como ferramenta de apoio, não substituindo a avaliação profissional. A responsabilidade clínica permanece com nutricionistas e demais especialistas, assegurando validade ética e jurídica quanto às recomendações fornecidas.

    \subsection{Métricas de Avaliação}

A avaliação do protótipo considera critérios de precisão e segurança das recomendações. Em particular, erros relacionados à sugestão de alimentos potencialmente perigosos para pacientes renais são tratados como criticidade máxima. Por essa razão, a métrica principal adotada é a \textbf{Taxa de Rejeição de Alimentos Nocivos (TRAN)}, definida como a proporção de alimentos clinicamente contraindicados que o sistema identifica e classifica corretamente como restritos ou proibidos. A escolha da TRAN como métrica central justifica-se pelo fato de que, em sistemas de recomendação aplicados à saúde, a consequência de um falso negativo (não rejeitar um alimento nocivo) é clinicamente mais grave do que a de um falso positivo (restringir desnecessariamente um alimento seguro) \cite{mckensy2022}. Também são analisados a \textbf{consistência das regras} aplicadas pelo raciocinador HermiT --- verificando se a ontologia permanece livre de contradições lógicas ---, a \textbf{aderência às diretrizes KDIGO/KDOQI} e a \textbf{clareza das justificativas semânticas} fornecidas pelo sistema, garantindo transparência e rastreabilidade no processo de recomendação.

% --- 5. ARQUITETURA DO SISTEMA ---
\section{Arquitetura do Sistema}

A arquitetura proposta segue uma abordagem orientada a conhecimento, na qual uma ontologia nutricional constitui o núcleo semântico do sistema. A Figura~\ref{fig:arquitetura} apresenta o fluxo completo, desde a entrada de dados do paciente até a geração das recomendações dietéticas.

\begin{figure*}[t]
\centering
\begin{tikzpicture}[
    node distance=1.2cm and 1.8cm,
    entrada/.style={rectangle, rounded corners, draw=black!70, fill=blue!10,
        text width=3cm, minimum height=1.1cm, align=center, font=\small\bfseries},
    processo/.style={rectangle, rounded corners, draw=black!70, fill=orange!10,
        text width=3cm, minimum height=1.1cm, align=center, font=\small\bfseries},
    ontologia/.style={cylinder, draw=black!70, fill=green!10,
        shape border rotate=90, aspect=0.3,
        text width=2.6cm, minimum height=1.3cm, align=center, font=\small\bfseries},
    saida/.style={rectangle, rounded corners, draw=black!70, fill=red!10,
        text width=3cm, minimum height=1.1cm, align=center, font=\small\bfseries},
    base/.style={cylinder, draw=black!70, fill=yellow!10,
        shape border rotate=90, aspect=0.3,
        text width=2.6cm, minimum height=1.1cm, align=center, font=\small\bfseries},
    seta/.style={-{Stealth[length=3mm]}, thick, black!70},
    rotulo/.style={font=\scriptsize\itshape, text=black!60},
]
\node[entrada] (paciente) {Interface Web\\(Flask)};
\node[processo, right=2cm of paciente] (classif) {Classificador\\de Estágio};
\node[ontologia, right=2cm of classif] (onto) {Ontologia\\OntoDRC\\(OWL 2)};
\node[processo, below=1.8cm of onto] (motor) {Motor de\\Inferência};
\node[base, left=2cm of motor] (taco) {Base TACO\\(Alimentos)};
\node[saida, below=1.8cm of motor] (saida) {Recomendação\\Dietética};
\node[saida, right=2cm of saida] (alerta) {Alertas\\Clínicos};
\node[entrada, left=2cm of saida] (nutri) {Avaliação do\\Nutricionista};
\draw[seta] (paciente) -- node[above, rotulo] {dados} (classif);
\draw[seta] (classif) -- node[above, rotulo] {estágio} (onto);
\draw[seta] (onto) -- node[right, rotulo] {restrições} (motor);
\draw[seta] (taco) -- node[above, rotulo] {composição} (motor);
\draw[seta] (motor) -- node[right, rotulo] {filtragem} (saida);
\draw[seta] (motor) -| node[above, rotulo, pos=0.3] {riscos} (alerta);
\draw[seta] (saida) -- node[above, rotulo] {apoio} (nutri);
\end{tikzpicture}
\caption{Arquitetura do Sistema de Recomendação Nutricional baseado em ontologias.}
\label{fig:arquitetura}
\end{figure*}

O sistema é composto por cinco componentes principais, descritos a seguir.

    \subsection{Interface de Entrada (Flask)}

O primeiro componente é a interface web, desenvolvida com o \textit{framework} Flask, responsável por coletar os dados do paciente. As informações solicitadas incluem: idade, peso corporal (em kg), Taxa de Filtração Glomerular (TFG), condição de diálise (sim/não) e presença de comorbidades como diabetes mellitus e hipertensão arterial. A TFG é o principal exame utilizado para avaliar a função renal: ela mede o volume de sangue filtrado pelos glomérulos renais por unidade de tempo, sendo expressa em mL/min/1,73m\textsuperscript{2}. Valores normais situam-se acima de 90 mL/min; à medida que a TFG diminui, a capacidade de filtragem dos rins se reduz, indicando progressão da DRC \cite{kdigo2024}. A interface é projetada para ser acessível e intuitiva, considerando que o público-alvo inclui pacientes atendidos pelo SUS com diferentes níveis de letramento digital.

    \subsection{Classificador de Estágio}

Com base na TFG informada, o sistema classifica automaticamente o paciente em um dos seis estágios da DRC conforme os critérios da KDIGO: G1 (TFG $\geq$ 90), G2 (60--89), G3a (45--59), G3b (30--44), G4 (15--29) e G5 ($<$ 15 mL/min/1,73m\textsuperscript{2}). Essa classificação determina quais restrições nutricionais serão aplicadas pelo motor de inferência. O classificador também considera se o paciente se encontra em hemodiálise ou diálise peritoneal, uma vez que pacientes em diálise possuem restrições e necessidades proteicas distintas dos demais estágios.

    \subsection{Ontologia OntoDRC}

O núcleo semântico do sistema é a ontologia OntoDRC, desenvolvida em OWL~2 e editada no Protégé~5.6. A ontologia organiza sete classes principais --- Alimento, Nutriente, Paciente, EstágioDRC, RestriçãoNutricional, RecomendaçãoDietética e RiscoClínico --- interligadas por propriedades semânticas e atributos quantitativos. A ontologia armazena tanto o conhecimento estrutural (relações entre conceitos) quanto os dados instanciados (alimentos da tabela TACO com seus valores nutricionais e limites de restrição por estágio).

As relações semânticas são formalizadas por meio de sete \textit{Object Properties}, cada uma com papel específico no raciocínio do sistema:

\begin{itemize}
    \item \texttt{contémNutriente} (Alimento $\rightarrow$ Nutriente): associa cada alimento aos nutrientes que contém em quantidade relevante. Por exemplo, a Banana contém Potássio.
    \item \texttt{possuiEstágio} (Paciente $\rightarrow$ EstágioDRC): vincula o paciente ao seu estágio atual da DRC, determinado pela TFG.
    \item \texttt{temRestrição} (EstágioDRC $\rightarrow$ RestriçãoNutricional): mapeia as restrições nutricionais aplicáveis a cada estágio, como o limite de 2.000~mg/dia de potássio para G3b.
    \item \texttt{recomendaAlimento} (RecomendaçãoDietética $\rightarrow$ Alimento): conecta uma recomendação gerada ao alimento sugerido como permitido.
    \item \texttt{restringeAlimento} (RestriçãoNutricional $\rightarrow$ Alimento): indica que determinada restrição se aplica a um alimento específico.
    \item \texttt{provocaRisco} (Nutriente $\rightarrow$ RiscoClínico): estabelece a relação causal entre o excesso de um nutriente e a complicação clínica associada, como Potássio $\rightarrow$ Hipercalemia.
    \item \texttt{geradaPara} (RecomendaçãoDietética $\rightarrow$ Paciente): vincula a recomendação ao paciente para o qual foi gerada, garantindo rastreabilidade.
\end{itemize}

Os atributos quantitativos --- como \texttt{potássioPor100g}, \texttt{sódioPor100g}, \texttt{temTFG} e \texttt{limiteDiárioMg} --- são representados como \textit{Data Properties}, possibilitando comparações numéricas durante a inferência.

    \subsection{Motor de Inferência}

O motor de inferência é implementado em Python~3 com a biblioteca \texttt{owlready2}, que permite carregar a ontologia OWL, consultar classes e indivíduos, e executar o raciocinador HermiT programaticamente. O processo de inferência segue três passos: (1)~identificação das restrições nutricionais associadas ao estágio do paciente via a relação \texttt{temRestrição}; (2)~comparação dos valores nutricionais de cada alimento cadastrado com os limites definidos nas restrições; e (3)~classificação de cada alimento como \textit{Permitido}, \textit{Restrito} ou \textit{Proibido}, acompanhada de uma justificativa semântica que explicita o nutriente excedente e o risco clínico associado.

As regras de classificação consideram os seguintes critérios principais, baseados nas diretrizes KDIGO e KDOQI:

\begin{itemize}
    \item \textbf{Potássio}: alimentos com mais de 200~mg/100g são classificados como restritos a partir do estágio G3b, e proibidos no estágio G5.
    \item \textbf{Sódio}: o limite diário recomendado é de 2.000~mg para todos os estágios a partir de G3a; alimentos ultraprocessados com alto teor de sódio são sinalizados.
    \item \textbf{Fósforo}: alimentos com mais de 150~mg/100g são restritos a partir de G3b; fontes de fósforo inorgânico (aditivos industriais) recebem alerta adicional.
    \item \textbf{Proteína}: a restrição varia conforme o estágio --- 0,6 a 0,8~g/kg/dia para estágios G3--G5 pré-diálise, e 1,0 a 1,2~g/kg/dia para pacientes em hemodiálise.
    \item \textbf{Líquidos}: restrição de 500 a 1.000~mL/dia acima do volume urinário residual para pacientes em diálise.
\end{itemize}

    \subsection{Saída: Recomendações e Alertas Clínicos}

O sistema gera dois tipos de saída. O primeiro é a \textbf{lista de recomendações}, que apresenta os alimentos avaliados com sua classificação (Permitido, Restrito ou Proibido) e a justificativa semântica correspondente --- por exemplo: ``Banana classificada como \textit{restrita}: potássio de 358~mg/100g excede o limite de 200~mg/porção para estágio G3b; risco associado: hipercalemia''. O segundo tipo é o \textbf{alerta clínico}, emitido automaticamente quando um alimento apresenta risco elevado em mais de um nutriente simultaneamente ou quando o perfil do paciente inclui comorbidades que intensificam o risco.

Todas as recomendações são acompanhadas de uma nota reforçando que a decisão final cabe ao nutricionista responsável, e que o sistema atua exclusivamente como ferramenta de apoio à decisão clínica.

    \subsection{Tecnologias Utilizadas}

A Tabela~\ref{tab:tecnologias} resume o \textit{stack} tecnológico adotado no desenvolvimento do protótipo.

\begin{table}[h]
\centering
\caption{Tecnologias e ferramentas do protótipo}
\label{tab:tecnologias}
\footnotesize
\begin{tabular}{|p{1.8cm}|p{2.5cm}|p{2.2cm}|}
\hline
\textbf{Componente} & \textbf{Tecnologia} & \textbf{Função} \\
\hline
Linguagem & Python 3.10+ & Implementação geral \\
\hline
Ontologia & OWL 2 / Protégé 5.6 & Modelagem semântica \\
\hline
Raciocinador & HermiT & Inferência \\
\hline
Manipulação & owlready2 & Acesso OWL \\
\hline
Interface & Flask & Entrada/saída \\
\hline
Base alimentar & TACO (Unicamp) & Composição \\
\hline
Diretrizes & KDIGO / KDOQI & Regras clínicas \\
\hline
\end{tabular}
\end{table}

% --- 6. RESULTADOS ESPERADOS ---
\section{Resultados Esperados}

Espera-se, como resultado central, a criação de um protótipo funcional capaz de oferecer orientações alimentares personalizadas para pacientes com Doença Renal Crônica. A ontologia servirá como base para organizar regras clínicas e representar relações entre alimentos, nutrientes e possíveis riscos. O sistema deve diferenciar itens permitidos, restritos e proibidos conforme o estágio da doença, apresentando justificativas claras e coerentes para cada recomendação. Além do desenvolvimento técnico, serão avaliadas a precisão e a robustez do modelo por meio de testes, validação das regras e métricas específicas, como a Taxa de Rejeição de Alimentos Nocivos (TRAN), bem como a clareza explicativa e a consistência semântica do mecanismo de inferência.

No âmbito social, o projeto pretende disponibilizar uma ferramenta acessível que diminua a dependência de conteúdos inadequados encontrados na internet, oferecendo apoio seguro e estruturado a pacientes atendidos pelo SUS, sem substituir o trabalho do nutricionista. Todos os resultados devem seguir rigorosamente os princípios da LGPD \cite{brasil2018}, assegurando anonimização, segurança da informação e tratamento ético dos dados. Ao final, espera-se entregar um protótipo confiável, fundamentado em diretrizes clínicas e alinhado às normas de privacidade na área da saúde.

% --- 7. CONSIDERAÇÕES FINAIS ---
\section{Considerações Finais}

Os resultados deste trabalho evidenciam que pacientes com Doença Renal Crônica enfrentam grandes dificuldades para encontrar informações nutricionais confiáveis, claras e adequadas ao seu quadro clínico. A predominância de recomendações genéricas e não validadas coloca em risco a segurança alimentar do paciente renal e agrava desigualdades de acesso ao conhecimento especializado. Nesse cenário, a integração entre Inteligência Artificial e ontologias mostra-se não apenas uma alternativa tecnicamente eficiente, mas uma necessidade estratégica para promover saúde pública de qualidade.

A arquitetura desenvolvida demonstra que a organização semântica do conhecimento nutricional, combinada ao uso de inferência automatizada, é capaz de gerar recomendações explicáveis, rastreáveis e aderentes às diretrizes clínicas. Esse modelo contribui para ampliar o acesso da população a informações seguras, apoiar profissionais do SUS e fortalecer práticas de cuidado baseadas em ciência. Todo o desenvolvimento foi guiado pelos princípios da Lei Geral de Proteção de Dados, assegurando sigilo, anonimização e governança ética no tratamento de informações sensíveis.

Conclui-se, portanto, que a união de IA e ontologias representa um avanço importante na democratização de informações nutricionais confiáveis para pacientes renais. O protótipo desenvolvido demonstra o potencial dessa abordagem para apoiar decisões clínicas, reduzir riscos alimentares e aproximar tecnologias emergentes das necessidades reais da população brasileira, reafirmando o compromisso com a equidade e a qualidade da saúde pública.

% --- REFERÊNCIAS ---
\begin{thebibliography}{99}
    \bibitem{posdoc}
    GOTTSCHALG-DUQUE, C. \textit{Proposta de desenvolvimento de um Sistema de Recomendação para a manutenção da Saúde de pacientes renais brasileiros}. Documento de Pós-Doutorado (manuscrito não publicado), Florida Tech, 2024.

    \bibitem{aggarwal}
    AGGARWAL, C. C. \textit{Recommender Systems: The Textbook}. Springer, 2016.

    \bibitem{kulkarni}
    KULKARNI, A. et al. \textit{Applied Recommender Systems with Python}. Apress, 2023.

    \bibitem{paper_onto}
    GOTTSCHALG-DUQUE, C. et al. Leveraging Ontologies for Intelligent Health Recommendation Systems. Manuscrito não publicado, 2023.

    \bibitem{felfernig}
    FELFERNIG, A. et al. \textit{Group Recommender Systems}. Springer, 2018.

    \bibitem{mendes2021}
    MENDES, A. et al. A semantic-based personalized food recommender system for older adults. \textit{Journal of Biomedical Informatics}, v. 118, 2021.

    \bibitem{zhang2023}
    ZHANG, X. et al. Ontology-driven nutrition recommendation for chronic disease management. \textit{Journal of Nutrition \& Health Sciences}, v. 12, n. 4, 2023.

    \bibitem{krieger2023}
    KRIEGER, C. S. A importância do Sistema Único de Saúde para o aprimoramento da saúde coletiva. \textit{RCMOS -- Revista Científica Multidisciplinar O Saber}, v. 3, n. 1, 2023.

    \bibitem{vasconcelos2021}
    VASCONCELOS, M. I. L. et al. Nutrição e doença renal crônica (DRC): Apresentação das novas recomendações e padrões alimentares conforme as últimas evidências científicas. \textit{Research, Society and Development}, v. 10, n. 6, e28610615891, 2021.

    \bibitem{gottschalg2025}
    GOTTSCHALG-DUQUE, C. et al. Leveraging Ontologies for Intelligent Health Recommendation Systems. In: \textit{ONTOBRAS}, 2025.

    \bibitem{ons}
    VITALI, F. et al. ONS: an ontology for a standardized description of interventions and observational studies in nutrition. \textit{Genes \& Nutrition}, v. 13, n. 1, art. 16, 2018. DOI: 10.1186/s12263-018-0601-y.

    \bibitem{ckdo}
    Chronic Kidney Disease Ontology (CKDO). Disponível em: \url{https://bioportal.bioontology.org/ontologies/CKDO}. Acesso em: 08 dez. 2025.

    \bibitem{mckensy2022}
    MCKENSY-SAMBOLA, D.; RODRÍGUEZ-GARCÍA, M. Á.; GARCÍA-SÁNCHEZ, F.; VALENCIA-GARCÍA, R. Ontology-based nutritional recommender system. \textit{Applied Sciences}, v. 12, n. 1, art. 143, 2022. DOI: 10.3390/app12010143.

    \bibitem{gogebakan2024}
    GÖĞEBAKAN, K. et al. A drug prescription recommendation system based on novel DIAKID ontology and extensive semantic rules. \textit{Health Information Science and Systems}, v. 12, n. 1, art. 27, 2024. DOI: 10.1007/s13755-024-00286-7.

    \bibitem{vitali2025}
    VITALI, F. et al. The Role and Applications of Semantic Interoperability Tools and eXplainable AI in the Development of Smart Food Systems. \textit{Computers and Electronics in Agriculture}, 2025.

    % NOVAS REFERÊNCIAS
    \bibitem{brasil2018}
    BRASIL. Lei n.º 13.709, de 14 de agosto de 2018. Lei Geral de Proteção de Dados Pessoais (LGPD). \textit{Diário Oficial da União}, Brasília, DF, 15 ago. 2018.

    \bibitem{botelho2021}
    BOTELHO, M. C.; CAMARGO, E. P. Proteção de dados pessoais de saúde na Lei Geral de Proteção de Dados. \textit{Revista de Direito Sanitário}, São Paulo, v. 21, e-0021, 2021.

    \bibitem{gruber1993}
    GRUBER, T. R. A translation approach to portable ontology specifications. \textit{Knowledge Acquisition}, v. 5, n. 2, p. 199--220, 1993.

    \bibitem{ibge2023}
    INSTITUTO BRASILEIRO DE GEOGRAFIA E ESTATÍSTICA (IBGE). Pesquisa Nacional por Amostra de Domicílios Contínua: Segurança Alimentar 2022--2023. Rio de Janeiro: IBGE, 2023.

    \bibitem{chen2019}
    CHEN, J. et al. The use of smartphone health apps and other mobile health (mHealth) technologies in dietetic practice: a three country study. \textit{Journal of Human Nutrition and Dietetics}, v. 30, n. 4, p. 439--452, 2019.

    \bibitem{kdigo2024}
    KIDNEY DISEASE: IMPROVING GLOBAL OUTCOMES (KDIGO). KDIGO 2024 Clinical Practice Guideline for the Evaluation and Management of Chronic Kidney Disease. \textit{Kidney International}, v. 105, n. 4S, p. S117--S314, 2024.

\end{thebibliography}
\end{document}
